This section moves beyond aggregate scores to analyze \emph{how} \cosa{} changes errors, \emph{where} gains occur at sample level, and \emph{why} behavior differs across datasets and backbones.

\subsection{Error Correction Mechanics: FP/FN Transitions}
To expose correction behavior, we compare baseline and \cosa{} error counts and report percentage reduction:
\begin{equation}
\begin{aligned}
\mathrm{FP\ Red.}&=\frac{\mathrm{FP}_{\text{base}}-\mathrm{FP}_{\text{cosa}}}{\mathrm{FP}_{\text{base}}}\times 100,\\
\mathrm{FN\ Red.}&=\frac{\mathrm{FN}_{\text{base}}-\mathrm{FN}_{\text{cosa}}}{\mathrm{FN}_{\text{base}}}\times 100.
\end{aligned}
\end{equation}
Positive values indicate fewer errors with \cosa{}; negative values indicate error increase.

\begin{table}[t]
  \centering
  \caption{Error correction mechanics when adding \cosa{}: percent reduction in FP/FN counts computed as $(\mathrm{baseline}-\mathrm{\cosa{}})/\mathrm{baseline}\times100$ (same test split, cached per-sample confusion). Positive values mean \cosa{} reduces that error count; negative values mean \cosa{} increases it.}
  \label{tab:error_mechanics}
  \setlength{\tabcolsep}{4pt}
  \renewcommand{\arraystretch}{1.05}
  \begin{tabular}{l l l}
    \toprule
    Model family & FP reduction & FN reduction \\
    \midrule
    CNN (FC-Siam) & -30.3\% & \textbf{+23.7\%} \\
    Attention (STANet) & \textbf{+68.5\%} & -179.8\% \\
    Transformer (BIT) & +13.6\% & -10.2\% \\
    \bottomrule
  \end{tabular}
\end{table}


\noindent\textbf{Interpretation.} Table~\ref{tab:error_mechanics} shows that the correction strategy is backbone-dependent. In the controlled CNN setting (FC-Siam), \cosa{} is recall-oriented: FN decreases while FP rises. In STANet, behavior flips to precision-oriented filtering: FP drops strongly while FN increases. In BIT, shifts are comparatively small, consistent with partial redundancy when temporal interaction is already strong \cite{chen2022bit,w2022,wei2024}. This confirms that \cosa{} is adaptive rather than uniformly precision- or recall-biased.

\subsection{Patch-Level Consistency and Sample-Wise Distribution}
For each sample $i$, we compute:
\begin{equation}
\Delta \mathrm{F1}_c^{(i)} = \mathrm{F1}_{c,\text{cosa}}^{(i)} - \mathrm{F1}_{c,\text{base}}^{(i)},
\end{equation}
then group samples into significant gain ($>10$), moderate gain ($2$--$10$), neutral ($-1$ to $+2$), and degradation ($<-1$) buckets.

\begin{table*}[t]
  \centering
  \caption{Patch-level consistency of \cosa{} improvements. Each cell reports the percentage of samples in a category and the average $\Delta$\Fonec{} (points) within that category.}
  \label{tab:consistency}
  \setlength{\tabcolsep}{10pt}
  \renewcommand{\arraystretch}{1.05}
  \small
  \begin{tabular}{l c c c}
    \toprule
    Category & FC-Siam & STANet & BIT \\
    \midrule
    Significant gain & \textbf{6.25\%}, +16.97 & 5.42\%, +19.13 & 2.10\%, \textbf{+90.50} \\
    Moderate gain & 21.09\%, +5.00 & \textbf{21.83\%}, \textbf{+5.52} & 1.42\%, +4.58 \\
    Neutral & 58.59\%, \textbf{+0.58} & 63.67\%, +0.05 & \textbf{89.31\%}, +0.04 \\
    Degradation & \textbf{14.06\%}, \textbf{-5.78} & 9.08\%, -22.97 & 7.18\%, -10.87 \\
    \bottomrule
  \end{tabular}
\end{table*}


\noindent\textbf{Interpretation.} Table~\ref{tab:consistency} shows that most samples are neutral across all families (especially BIT), indicating that refinement is largely non-interfering. For FC-Siam and STANet, improvement buckets (moderate + significant) exceed degradation frequency, which explains positive aggregate gains. BIT remains predominantly neutral, consistent with its near-ceiling baseline and marginal overall gain in standalone tests.

\noindent\textbf{When \cosa{} helps.} Strongest improvements are concentrated in FN-heavy and boundary-ambiguous samples, where local correlation cues recover coherent changed regions missed by baseline predictions.

\noindent\textbf{When \cosa{} can hurt.} Degradation is concentrated in ambiguous or noisy-label regions and weak-texture areas, where correlation structure is less reliable.

\subsection{Cross-Dataset Adaptive Behavior}
Controlled results in Table~\ref{tab:fcsiam_vs_cosa} show that \cosa{} improves \Fonec{} on all four datasets, but via different precision/recall routes. LEVIR-CD benefits mainly from recall recovery with stable precision; S2Looking improves both precision and recall; DSIFN shifts strongly toward recall; and CLCD shifts toward precision filtering under noisier boundaries. This pattern supports the method claim from Section~\ref{sec:method}: learnable residual gating enables dataset-dependent calibration rather than fixed global behavior.

\subsection{Computational Cost and Efficiency Discussion}
\cosa{} is designed as a local decoder refinement with neighborhood-based correlation sampling and lightweight residual gating. For window size $k$ and channel count $C$, the dominant cost is local correlation sampling $O(HWk^2C)$ plus linear-time aggregation/gating $O(HWC)$, which avoids global all-pairs interaction.

In our controlled FC-Siam setup, the parameter increase is negligible (\(\Delta\)Params \(= +0.000066\)M; see repository efficiency summary), while delivering +1.75 \Fonec{} on LEVIR-CD. This profile supports practical deployment as a plug-in module. A standardized cross-hardware latency benchmark is left for future work.
